\subsection{The map \texorpdfstring{\(p\)}{p}}

\begin{lemma}\label{sga3:XXI:1.2.1}
Let \(p:M\to M^\ast\) (respectively \(V\to V^\ast\)) be the linear map
defined by
\begin{equation}
 p(x)=\sum_{u\in R^\ast}(u,x)u.
 \tag{6}\label{sga3:XXI:eq:6}
\end{equation}
Write \(\ell(x)=(p(x),x)\). Then:
\begin{align}
 \ell(x)&\geq0,\qquad \ell(\alpha)>0
     &&\text{for }\alpha\in R,
 \tag{7}\label{sga3:XXI:eq:7}\\
 \ell(wx)&=\ell(x)
     &&\text{for }w\in W,
 \tag{8}\label{sga3:XXI:eq:8}\\
 (p(x),y)&=(p(y),x)
     &&\text{for }x,y\in V,
 \tag{9}\label{sga3:XXI:eq:9}\\
 \ell(\alpha)\alpha^\ast&=2p(\alpha)
     &&\text{for }\alpha\in R.
 \tag{10}\label{sga3:XXI:eq:10}
\end{align}
\end{lemma}

\begin{proof}
The first three relations are immediate.
We prove the last one. By \eqref{sga3:XXI:eq:1}, for \(u\in V^\ast\),
\begin{align*}
 (u,\alpha)^2\alpha^\ast
  &=(u,\alpha)u-(u,\alpha)s_\alpha(u)\\
  &=(u,\alpha)u+(u,s_\alpha(\alpha))s_\alpha(u)\\
  &=(u,\alpha)u+(s_\alpha(u),\alpha)s_\alpha(u)
       \qquad\bigl(s_\alpha^2=\operatorname{id}\bigr).
\end{align*}
Since \(s_\alpha\) permutes \(R^\ast\) by
\((\mathrm{DR\,II})\), summing over \(u\in R^\ast\) gives the result.
\end{proof}

\begin{scholium}\label{sga3:XXI:1.2.2}
Relation \eqref{sga3:XXI:eq:10} says that the map
\(\alpha\mapsto\ell(\alpha)\alpha^\ast\) is the restriction to \(R\)
of a linear map from \(M\) to \(M^\ast\). In particular,
\[
 (-\alpha)^\ast=-\alpha^\ast.
\]
\end{scholium}

\begin{corollary}\label{sga3:XXI:1.2.3}
The map \(p\) induces an isomorphism from \(\VR(R)\) onto
\(\VR(R^\ast)\).
\end{corollary}

\begin{proof}
The map \(p\) sends \(\VR(R)\) onto \(\VR(R^\ast)\). Hence
\[
 \dim\VR(R)\geq\dim\VR(R^\ast).
\]
Applying this inequality to the dual root datum gives
\[
 \dim\VR(R)=\dim\VR(R^\ast),
\]
so the surjective map \(p\) is also bijective.
\end{proof}

\begin{corollary}\label{sga3:XXI:1.2.4}
One has \(\dim\VR(R)=\dim\VR(R^\ast)\), and therefore
\[
 \rgss(\cR)=\rgss(\cR^\ast).
\]
\end{corollary}

\begin{corollary}\label{sga3:XXI:1.2.5}
The bilinear form \((u,x)\) is nondegenerate on
\(\VR(R^\ast)\times\VR(R)\), and thus puts these
\(\mathbb Q\)-vector spaces in duality.
\end{corollary}

\begin{proof}
If \((u,x)=0\) for every \(u\in R^\ast\), then \(p(x)=0\).
\end{proof}

\begin{corollary}\label{sga3:XXI:1.2.6}
The symmetric bilinear form \((p(x),y)\) is positive and nondegenerate
on \(\VR(R)\).
\end{corollary}

\begin{corollary}\label{sga3:XXI:1.2.7}
The group \(W\) acts faithfully on \(R\) (and hence on the other sets
listed in Scholium~\ref{sga3:XXI:1.1.13}).
\end{corollary}

\begin{proof}
Let \(u\in V^\ast\), and suppose \(w\in W\) satisfies
\(w(\alpha)=\alpha\) for every \(\alpha\in R\). We prove \(w(u)=u\).
For every \(\alpha\in R\),
\[
 (w(u)-u,\alpha)
 =(w(u),\alpha)-(u,\alpha)
 =(u,w^{-1}(\alpha))-(u,\alpha)=0.
\]
But \(w(u)-u\in\VR(R^\ast)\). Since it is orthogonal to every root, it
vanishes by Corollary~\ref{sga3:XXI:1.2.5}.
\end{proof}

\begin{corollary}\label{sga3:XXI:1.2.8}
The group \(W\) is finite.
\end{corollary}

\begin{proposition}\label{sga3:XXI:1.2.9}
The operations of \(W\) respect the correspondence between roots and
coroots. In other words, for \(\alpha\in R\) and \(w\in W\),
\[
 w(\alpha^\ast)=w(\alpha)^\ast.
\]
\end{proposition}

\begin{proof}
It suffices to verify this for \(w=s_\beta\), \(\beta\in R\), that is,
to verify
\[
 s_\beta(\alpha^\ast)=s_\beta(\alpha)^\ast.
\]
Now \(\ell(s_\beta(\alpha))s_\beta(\alpha)^\ast/2\) equals
\begin{align*}
 p(s_\beta(\alpha))
 &=\sum_{u\in R^\ast}(u,s_\beta(\alpha))u\\
 &=\sum_{u\in R^\ast}(s_\beta(u),\alpha)u\\
 &=\sum_{u\in R^\ast}(u,\alpha)s_\beta(u)
  =s_\beta(p(\alpha)).
\end{align*}
Since \(\ell(s_\beta(\alpha))=\ell(\alpha)\), we obtain
\[
 s_\beta(\alpha)^\ast
 =s_\beta\!\left(\frac{2p(\alpha)}{\ell(\alpha)}\right)
 =s_\beta(\alpha^\ast).
\]
\end{proof}

\begin{corollary}\label{sga3:XXI:1.2.10}
If \(\alpha\in R\) and \(w\in W\), then
\[
 ws_\alpha w^{-1}=s_{w(\alpha)}.
\]
\end{corollary}

\begin{proof}
For \(x\in V\),
\[
 ws_\alpha w^{-1}(x)
 =x-(\alpha^\ast,w^{-1}(x))w(\alpha)
 =x-(w(\alpha^\ast),x)w(\alpha).
\]
By Proposition~\ref{sga3:XXI:1.2.9}, the last expression is
\[
 \begin{aligned}
 x-(w(\alpha)^\ast,x)w(\alpha)
   &=s_{w(\alpha)}(x).
 \end{aligned}
\]
\end{proof}

\begin{corollary}\label{sga3:XXI:1.2.11}\leavevmode\par
Let \(R'\subset R\). The quadruple
\((M,M^\ast,R',R'^\ast)\) is a root datum if and only if
\(\alpha,\beta\in R'\) implies \(s_\alpha(\beta)\in R'\).
\end{corollary}

\section{Relations between two roots}

\subsection{Proportional roots}

\begin{proposition}\label{sga3:XXI:2.1.1}
Let \(\alpha\) and \(\beta\) be two roots. The following conditions are
equivalent:
\begin{enumerate}
\item[(i)] there exists \(k\in\mathbb Q\) such that
  \(\alpha=k\beta\);
\item[(ii)] \(s_\alpha=s_\beta\).
\end{enumerate}
Under these conditions,
\(\alpha^\ast=k^{-1}\beta^\ast\), and \(k\) is one of
\[
 1,\ -1,\ 2,\ -2,\ \tfrac12,\ -\tfrac12.
\]
\end{proposition}

\begin{proof}
Assume first \SGARefLink{sga3:XXI:2.1.1:item:i:02}{(i)}. Then
\[
 \alpha^\ast
 =\ell(\alpha)^{-1}2p(\alpha)
 =k^{-2}\ell(\beta)^{-1}2k\,p(\beta)
 =k^{-1}\beta^\ast.
\]
This immediately gives \(s_\alpha=s_\beta\). Conversely, if
\(s_\alpha=s_\beta\), then
\[
 \alpha=s_\alpha(-\alpha)
       =s_\beta(-\alpha)
       =(\beta^\ast,\alpha)\beta-\alpha,
\]
so
\[
 2\alpha=(\beta^\ast,\alpha)\beta,
\qquad(\beta^\ast,\alpha)\in\mathbb Z;
\]
thus \SGARefLink{sga3:XXI:2.1.1:item:ii:02}{(ii)} implies \SGARefLink{sga3:XXI:2.1.1:item:i:02}{(i)}. Finally, if \(\alpha=k\beta\), then
\(\alpha^\ast=k^{-1}\beta^\ast\), whence
\[
 (\alpha^\ast,\beta)=2k^{-1},
 \qquad(\beta^\ast,\alpha)=2k.
\]
Thus \(2k\) and \(2k^{-1}\) are integers, which gives the stated list.
\end{proof}

\begin{application}\label{sga3:XXI:2.1.2}
Root data \(\cR\) with \(\rgss(\cR)=1\) are of one of the following two
types:
\begin{enumerate}
\item[\SGAInlineRefTarget{sga3:XXI:2.1.1:item:i:02}(i)] Type \(A_1\): there exists \(\alpha\in M\) such that the
  roots are \(\alpha\) and \(-\alpha\). The coroots are then
  \(\alpha^\ast\) and \(-\alpha^\ast\).
\item[\SGAInlineRefTarget{sga3:XXI:2.1.1:item:ii:02}(ii)] Type \(A'_1\): there exists \(\alpha\in M\) such that the
  roots are
  \(\alpha,-\alpha,2\alpha,-2\alpha\). The coroots are then
  \(\alpha^\ast,-\alpha^\ast,\alpha^\ast/2,-\alpha^\ast/2\).
\end{enumerate}
\end{application}
